\documentclass[10pt,conference]{IEEEtran}
\usepackage{cite}
\usepackage{amsmath,amssymb,amsfonts}
\usepackage{algorithmic}
\usepackage{graphicx}
\usepackage{textcomp}
\usepackage{xcolor}
\usepackage[hyphens]{url}
\usepackage{fancyhdr}
\usepackage{hyperref}


% Ensure letter paper
\pdfpagewidth=8.5in
\pdfpageheight=11in

\newcommand{\hpcayear}{2027}


%%%%%%%%%%%%%%%%%%%%%%%%%%%%%%%%%%%%%%%%
%%%%%%%%%%%%%% -- UPDATE -- %%%%%%%%%%%%%%%
\newcommand{\hpcasubmissionnumber}{NaN}
\title{Guidelines for Submission to HPCA \hpcayear{}}
%%%%%%%%%%%%%%%%%%%%%%%%%%%%%%%%%%%%%%%%


%%%%%%%%%%%%%%%%%%%%%%%%%%%%%%%%%%%%%%%%
%%%%%%%% -- ONLY FOR CAMERA READY -- %%%%%%%%
%\def\hpcacameraready{} % Uncomment to build camera-ready version
\newcommand{\hpcapubid}{0000--0000/00\$00.00}
\newcommand\hpcaauthors{First Author$\dagger$ and Second Author$\ddagger$}
\newcommand\hpcaaffiliation{First Affiliation$\dagger$, Second Affiliation$\ddagger$}
\newcommand\hpcaemail{Email(s)}

%%%%% -- ARTEFACT EVALUATION RESULTS -- %%%%%%
% Uncomment the following based on the badges that were awarded to this paper
%\def\aeopen{}           % The artifact is publically available
%\def\aereviewed{}     % The artefact has been reviewed
%\def\aereproduced{} % The results have been reproduced
%%%%%%%%%%%%%%%%%%%%%%%%%%%%%%%%%%%%%%%%

%%%%%%%%%%%%%%%%%%%%%%%%%%%%%%%%%%%%%
%%%%%%%%%% -- DO NOT MODIFY -- %%%%%%%%%%
%%%%%%%%%%%%%%%%%%%%%%%%%%%%%%%%%%%%%

\author{
  \ifdefined\hpcacameraready
    \IEEEauthorblockN{\hpcaauthors{}}
      \IEEEauthorblockA{
        \hpcaaffiliation{} \\
        \hpcaemail{}
      }
  \else
    \IEEEauthorblockN{\normalsize{HPCA \hpcayear{} Submission
      \textbf{\#\hpcasubmissionnumber{}}} \\
      \IEEEauthorblockA{
        Confidential Draft \\
        Do NOT Distribute!!
      }
    }
  \fi 
}

% Heading and footer for title page
\fancypagestyle{camerareadyfirstpage}{%
  \fancyhead{}
  \renewcommand{\headrulewidth}{0pt}
  \fancyhead[C]{
    \ifdefined\aeopen
    \parbox[][12mm][t]{13.5cm}{\hpcayear{} IEEE International Symposium on High-Performance Computer Architecture (HPCA)}    
    \else
      \ifdefined\aereviewed
      \parbox[][12mm][t]{13.5cm}{\hpcayear{} IEEE International Symposium on High-Performance Computer Architecture (HPCA)}
      \else
      \ifdefined\aereproduced
      \parbox[][12mm][t]{13.5cm}{\hpcayear{} IEEE International Symposium on High-Performance Computer Architecture (HPCA)}
      \else
      \parbox[][0mm][t]{13.5cm}{\hpcayear{} IEEE International Symposium on High-Performance Computer Architecture (HPCA)}
    \fi 
    \fi 
    \fi 
    \ifdefined\aeopen 
      \includegraphics[width=12mm,height=12mm]{ae-badges/open-research-objects.pdf}
    \fi 
    \ifdefined\aereviewed
      \includegraphics[width=12mm,height=12mm]{ae-badges/research-objects-reviewed.pdf}
    \fi 
    \ifdefined\aereproduced
      \includegraphics[width=12mm,height=12mm]{ae-badges/results-reproduced.pdf}
    \fi
  }
  %\fancyfoot[L]{\hpcapubid{} \copyright \hpcayear{} IEEE}
  \fancyfoot[C]{}
}
% Heading and footer for remaining pages
\fancyhead{}
\renewcommand{\headrulewidth}{0pt}
%\fancyhead[C]{\hpcayear{} IEEE International Symposium on
% High-Performance Computer Architecture (HPCA)}


\begin{document}
\maketitle

%Enables the camera ready header and footer
\ifdefined\hpcacameraready 
  \thispagestyle{camerareadyfirstpage}
  \pagestyle{empty}
\else
  \thispagestyle{plain}
  \pagestyle{plain}
\fi

\newcommand{\hpcaheight}{0mm}
\ifdefined\eaopen
\renewcommand{\hpcaheight}{12mm}
\fi
  


%%%%%%%%%%%%%%%%%%%%%%%%%%%%%%%%%%%%%%%%
%%%%%%%% -- PAPER CONTENT STARTS -- %%%%%%%%%

\begin{abstract}

  This document is intended to serve as the acceptable formatting
  sample for submissions to the IEEE International Symposium on
  High-Performance Computer Architecture (HPCA \hpcayear{}). This
  document provides guidelines that authors should follow when
  submitting papers to the conference.  This format is derived from
  the IEEE conference template IEEEtran.cls file with the objective of
  keeping the submission similar to the final version, i.e., the
  IEEEtran.cls template will also be used for the camera-ready
  version.

\end{abstract}

\section{Introduction}

This document provides instructions for submitting papers to HPCA
\hpcayear{}.  In an effort to respect the efforts of reviewers and in
the interest of fairness to all prospective authors, we request that
all submissions to HPCA \hpcayear{} follow the formatting and
submission rules detailed below. Submissions that violate these
instructions may not be reviewed, at the discretion of the Program
Chair, in order to maintain a review process that is fair to all
potential authors. This document is itself formatted using the HPCA
\hpcayear{} submission format. The content of this document mirrors
that of the submission instructions that appear on the conference
website. All questions regarding paper formatting and submission
should be directed to the Program Chair.


\subsection{Format Highlights}

Here are the format highlights in a nutshell:
\begin{itemize}
\item Paper must be submitted in printable PDF format.
\item Text must be in a minimum 10pt Times font, see Table~\ref{table:formatting}.
\item Papers must be at most 11 pages (not including references) in a
  two-column format.
\item No page limit for references.
\item Each reference must specify {\em all} authors, i.e., no {\em et al.}
\end{itemize}

\subsection{Paper Evaluation Objectives} 
The committee will make every effort to judge each submitted paper on
its own merits. There will be no target acceptance rate. We expect to
accept a wide range of papers with appropriate expectations for
evaluation --- while papers that build on significant past work with
strong evaluations are valuable, papers that open new areas with less
rigorous evaluation are equally welcome and especially encouraged. We
also acknowledge the wide range of evaluation methodologies including
modeling, simulation, prototyping, experimental implementation, real
product evaluation, etc.

\section{Paper Preparation Instructions}

\subsection{Paper Formatting}
Papers must be submitted in printable PDF format and should contain a
{\em maximum of 11 pages} of single-spaced two-column text, {\bf not
  including references}.  You may include any number of pages for
references, but see below for more instructions.  If you are using
\LaTeX~\cite{lamport94} to typeset your paper, then you should use the
template used to prepare this document, which you can find on the HPCA
\hpcayear{} website. If you use a different software package to
typeset your paper, then you should adhere to the guidelines given in
Table~\ref{table:formatting}.

\begin{scriptsize}
\begin{table}[h!]
  \centering
  \caption{Formatting guidelines for submission.}
  \label{table:formatting}
  \begin{tabular}{|l|l|}
    \hline
    \textbf{Field} & \textbf{Value}\\
    \hline
    \hline
    File format & PDF \\
    \hline
    Page limit & 11 pages, \textbf{not including references} \\
    \hline
    Paper size & US Letter 8.5in $\times$ 11in\\
    \hline
    Top margin & 0.7in\\
    \hline
    Bottom margin & 1in\\
    \hline
    Left margin & 0.7in\\
    \hline
    Right margin & 0.7in\\
    \hline
    Body & 2-column, single-spaced\\
    \hline
    Space between columns & 0.1in\\
    \hline
    Line spacing (leading) & 12pt \\
    \hline
    Body font & 10pt, Times\\
    \hline
    Abstract font & \textbf{9pt, Times, bold} \\
    \hline
    Section heading font & \textsc{12pt, Small Caps,} \\
	                 & \textsc{Roman Numbering (I, II, ...)} \\
    \hline
    Subsection heading font & \textit{10pt, italic, Alpha numering (A, B, ...)} \\
    \hline
    Caption font & \footnotesize{9pt (minimum)} \\
    \hline
    References & \scriptsize{8pt, no page limit, list all authors' names} \\
    \hline
  \end{tabular}
\end{table}
\end{scriptsize}

{\em Please ensure that you include page numbers with your
  submission}. This makes it easier for the reviewers to refer to
different parts of your paper when they provide comments. Please
ensure that your submission has a banner at the top of the title page,
similar to this document, which contains the submission number and the
notice of confidentiality.  If using the template, just replace `NaN'
with your submission number. 

\subsection{Content}
Reviewing will be {\em double blind} (no author list); therefore,
please do not include any author names on any submitted documents
except in the space provided on the submission form.  You must also
ensure that the metadata included in the PDF does not give away the
authors. If you are improving upon your prior work, refer to your
prior work in the third person and include a full citation for the
work in the bibliography.  For example, if you are building on {\em
  your own} prior work in the
papers~\cite{nicepaper1,nicepaper2,nicepaper3}, you would say
something like: ``While the authors of~\cite{nicepaper1,nicepaper2,nicepaper3} did X,
Y, and Z, this paper additionally does W, and is therefore much
better.''  Do NOT omit or anonymize references for blind review.  There
is one exception to this for your own prior work that appeared in IEEE
CAL, arXiv, workshops without archived proceedings, etc.\ as
discussed later in this document. 

\noindent\textbf{Figures and Tables:} Ensure that the figures and
tables are legible.  Please also ensure that you refer to your figures
in the main text.  Many reviewers print the papers in
gray-scale. Therefore, if you use colors for your figures, ensure that
the different colors are highly distinguishable in gray-scale. 

\noindent\textbf{References:}  There is no length limit for
references. {\em Each reference must explicitly list all authors of
  the paper.  Papers not meeting this requirement will be rejected.}
Since there is no length limit for the number of pages
used for references, there is no need to save space here. 


\section{Paper Submission Instructions}

\subsection{Guidelines for Determining Authorship}
IEEE guidelines dictate that authorship should be based on a {\em
  substantial intellectual contribution}. It is assumed that all
authors have had a significant role in the creation of an article that
bears their names. In particular, the authorship credit must be
reserved only for individuals who have met each of the following
conditions: 

\begin{enumerate}
\item Made a significant intellectual contribution to the theoretical
  development, system or experimental design, prototype development,
  and/or the analysis and interpretation of data associated with the
  work contained in the article; 

\item Contributed to drafting the article or reviewing and/or revising
  it for intellectual content; and 

\item Approved the final version of the article as accepted for
  publication, including references.
\end{enumerate}

A detailed description of the IEEE authorship guidelines and
responsibilities is available
online.\footnote{\url{https://www.ieee.org/publications_standards/publications/rights/Section821.html}}
Please keep these guidelines in mind while determining the author list
of your paper.

\subsection{Declaring Authors}
Declare all the authors of the paper upfront. Addition/removal of
authors once the paper is accepted will have to be approved by the
Program Chair, since it potentially undermines the goal of eliminating
conflicts for reviewer assignment.

\subsection{Areas and Topics}
Authors should indicate specific topics covered by the paper on the
submission page. If you are unsure whether your paper falls within the
scope of the conference, please check with the Program Chair --- HPCA
is a broad, multidisciplinary conference and encourages new topics.

\subsection{Declaring Conflicts of Interest}

Authors must register all their conflicts on the paper submission
site. Conflicts are needed to ensure appropriate assignment of
reviewers. If a paper is found to have an undeclared conflict that
causes a problem OR if a paper is found to declare false conflicts in
order to abuse or ``game'' the review system, the paper may be
rejected.

Please declare a conflict of interest with the following people for
any author of your paper.  A conflict occurs in the following cases:
\begin{enumerate}
\item Between advisors and advisees, forever. 
\item Between family members, forever (if they might be potential reviewers). 
\item Between people who have collaborated in the last 4 years. This
  collaboration can consist of a joint research or development
  project, a joint paper, or when there is direct funding from the
  potential reviewer (as opposed to company funding) to an author of
  the paper. Co-participation in professional activities, such as
  tutorials or studies, is not a cause for conflict. When in doubt,
  the author should check with the Program Chair.
\item Between people from the same institution, who were in the same
  institution in the last 4 years, or where one of them is being
  actively considered for employment (e.g. invitation for interview)
  by the other person's institution.
\item When there is funding involving the potential reviewer (as
  opposed to merely funding from the reviewer's company).
\item Between people who are members of the same research center where
  unpublished work is regularly discussed (e.g., at review meetings).
\item Between people whose relationship prevents the reviewer from
  being objective in their assessment.
\item There may be others not covered by the above with whom you
  believe a COI exists, for example, an ongoing collaboration which
  has not yet resulted in the creation of a paper or proposal. Please
  report such COIs; however, you may be asked to justify them. Please
  be reasonable. For example, you cannot declare a COI with a reviewer
  just because that reviewer works on topics similar to or related to
  those in your paper. The Program Chair may contact co-authors to
  explain a COI whose origin is unclear.
\end{enumerate}
The following scenarios \textit{do not} constitute a conflict:
\begin{enumerate}
\item Authors of previously-published, closely related work on that
  basis alone.
\item ``Service'' collaborations such as co-authoring a report for a
  professional organization, serving on a program committee, or
  co-presenting tutorials.
\item Co-authoring a paper that is a compendium of various projects,
  community-wide tools (e.g., gem5), non-research articles, and
  working groups (e.g. RISC-V) with no true research collaboration
  among the co-authors.
\item People who work on topics similar to or related to those in your
  papers.
\item People under the same umbrella funding award where there is no
  close collaboration, no discussion of unpublished work, and no joint
  benefit in the paper being published.
\end{enumerate}

Most reviews will be solicited among the members of the PC, but other
members from the community may also write reviews. Please declare all
your conflicts (not just restricted to the PC) on the submission
form. When in doubt, contact the Program Chair.

\subsection{ArXiv Submissions}
Authors may submit their work to HPCA \hpcayear{} for publication even
if the paper has been submitted to or currently appears on
Arxiv. However, please do recognize that Arxiv submissions cause
serious issues with the double-blind review process. The PC chair has
already advised the PC members not to consider Arxiv submissions in
their evaluation. However, we encourage the authors to reduce the
possibility of their name being disclosed. While there is no magical
solution to hide from the power of web search in locating Arxiv
papers, we encourage authors to avoid substantially similar titles
abstract etc., to reduce this match.

\subsection{Concurrent Submissions and Workshops}
By submitting a manuscript to HPCA \hpcayear{}, the authors guarantee
that the manuscript has not been previously published or accepted for
publication in a substantially similar form in any conference,
journal, or the archived proceedings of a workshop (e.g., in the
ACM/IEEE digital libraries) --- see exceptions below. The authors also
guarantee that no paper that contains significant overlap with the
contributions of the submitted paper will be under review for any
other conference or journal or an archived proceedings of a workshop
during the HPCA \hpcayear{} review period. Violation of any of these
conditions will lead to rejection.

The only exceptions to the above rules are for the authors' own papers
in (1) workshops without archived proceedings such as in the ACM/IEEE
digital libraries (or where the authors chose not to have their paper
appear in the archived proceedings), or (2) venues such as IEEE CAL or
arXiv where there is an explicit policy that such publication does not
preclude longer conference submissions.  In all such cases, the
submitted manuscript may ignore the above work to preserve author
anonymity. This information must, however, be provided on the
submission form --- the Program Chair will make this information
available to reviewers if it becomes necessary to ensure a fair
review.  As always, if you are in doubt, it is best to contact Program
Chair.


Finally, the ACM/IEEE Plagiarism Policies\footnote{\url{http://www.acm.org/publications/policies/plagiarism_policy}\\
  \url{https://www.ieee.org/publications_standards/publications/rights/plagiarism_FAQ.html}}
cover a range of ethical issues concerning the misrepresentation of
other works or one's own work.

\subsection{Guidelines on the use of AI-Generated Content}

In accordance with IEEE author guidelines and
policies~\footnote{\url{https://conferences.ieeeauthorcenter.ieee.org/author-ethics/guidelines-and-policies/submission-policies/}},
the use of content generated by artificial intelligence (AI) in a
paper submission (including but not limited to text, figures, images,
and code) shall be disclosed in the acknowledgments section of any
article submitted to an IEEE publication. The AI system used shall be
identified, and specific sections of the article that use AI-generated
content shall be identified and accompanied by a brief explanation
regarding the level at which the AI system was used to generate the
content.

To this end, papers with AI-Generated content must include an Appendix
after the references, and not counted in the page count, which
provides this information.

The use of AI systems for editing and grammar enhancement is common
practice and, as such, is generally outside the intent of the above
policy. In this case, disclosure as noted above is not required, but
recommended.


\section*{Acknowledgements}
This document is an updated version of HPCA 2022, 2023 and 2025 which,
in turn, has been derived from two previous conferences, in particular
HPCA 2021 and MICRO 2021, which, in turn, are derived from past MICRO,
HPCA, ISCA, and ASPLOS conferences.

%%%%%%% -- PAPER CONTENT ENDS -- %%%%%%%%


%%%%%%%%% -- BIB STYLE AND FILE -- %%%%%%%%
\bibliographystyle{IEEEtranS}
\bibliography{refs}
%%%%%%%%%%%%%%%%%%%%%%%%%%%%%%%%%%%%

\end{document}

